\documentclass[preview]{standalone}

\usepackage{amsmath}
\usepackage{amssymb}
\usepackage{stellar}
\usepackage{definitions}

\begin{document}

\id{sylow-theorems-applications}
\genpage

\section{Normal Sylow subgroups}

\begin{snippetcorollary}{sylow-subgroup-normal-iff-unique}{Normal Sylow subgroups}
    Let \(G\) be a finite \group. A Sylow \(p\)-subgroup of \(G\) is normal
    \ifandonlyif it is the unique Sylow \(p\)-subgroup of \(G\).
\end{snippetcorollary}

\begin{snippetproof}{sylow-subgroup-normal-iff-unique-proof}{sylow-subgroup-normal-iff-unique}{Normal Sylow subgroups}
    Sylow \(p\)-subgroups are all conjugate. A normal subgroup has no distinct
    conjugates, while a unique subgroup of a given order is necessarily fixed
    by conjugation.
\end{snippetproof}

\begin{snippetproposition}{sylow-normalizer-self-normalizing}{Normalizer of the normalizer of a Sylow subgroup}
    Let \(P\) be a Sylow \(p\)-subgroup of a finite \group \(G\). Then
    \[
        \groupnormalizer_G(\groupnormalizer_G(P))=\groupnormalizer_G(P).
    \]
\end{snippetproposition}

\begin{snippetproof}{sylow-normalizer-self-normalizing-proof}{sylow-normalizer-self-normalizing}{Normalizer of the normalizer of a Sylow subgroup}
    Let \(K=\groupnormalizer_G(P)\). Then \(P\normalsubgrp K\), and \(P\) is
    a Sylow \(p\)-subgroup of \(K\), so it is the unique such subgroup of
    \(K\). If \(g\in\groupnormalizer_G(K)\), then \(P^g\) is also a Sylow
    \(p\)-subgroup of \(K^g=K\), hence \(P^g=P\). Thus
    \(g\in\groupnormalizer_G(P)=K\). The reverse inclusion
    \(K\subgroupleq\groupnormalizer_G(K)\) always holds.
\end{snippetproof}

\begin{snippetproposition}{product-normal-subgroups-order}{Normal subgroups of pairwise coprime orders}
    Let \(G\) be a finite \group and let
    \(H_1,\ldots,H_n\normalsubgrp G\) have pairwise coprime orders. Their
    product is an internal direct product and
    \[
        \left|\prod_{i=1}^nH_i\right|=\prod_{i=1}^n|H_i|.
    \]
\end{snippetproposition}

\begin{snippetproof}{product-normal-subgroups-order-proof}{product-normal-subgroups-order}{Normal subgroups of pairwise coprime orders}
    Induct on \(n\). By induction,
    \(K=H_1\cdots H_{n-1}\) is normal and has order
    \(\prod_{i<n}|H_i|\). Thus \(|K|\) and \(|H_n|\) are coprime, so
    \(K\intersection H_n=1\) by Lagrange's theorem and
    \(|KH_n|=|K||H_n|\). Repeating the argument after isolating any factor
    shows
    \[
        H_i\intersection\prod_{j\neq i}H_j=1.
    \]
    These are precisely the internal direct-product conditions.
\end{snippetproof}

\begin{snippettheorem}{normal-sylow-direct-product-decomposition}{Decomposition by normal Sylow subgroups}
    If every Sylow subgroup of a finite \group \(G\) is normal, then \(G\) is
    their internal direct product. In particular, every finite
    \abeliangroup is the direct product of its Sylow subgroups.
\end{snippettheorem}

\begin{snippetproof}{normal-sylow-direct-product-decomposition-proof}{normal-sylow-direct-product-decomposition}{Decomposition by normal Sylow subgroups}
    Choose one Sylow \(p\)-subgroup \(P_p\) for each prime dividing \(|G|\).
    Their orders are pairwise coprime, and their product has order
    \(\prod_p|P_p|=|G|\) by the preceding proposition. Hence their product is
    all of \(G\) and is direct. In an abelian group every subgroup is normal.
\end{snippetproof}

\section{Groups of order pq}

\begin{snippettheorem}{p-q-groups-distinct-primes-product-order-theorem}{Classification of groups of order \(pq\)}
    Let \(p<q\) be distinct primes.
    \begin{enumerate}
        \item If \(q\not\equiv1\pmod p\), there is one group of order \(pq\)
            up to isomorphism: the cyclic group \(C_{pq}\).
        \item If \(q\equiv1\pmod p\), there are two groups of order \(pq\)
            up to isomorphism: \(C_{pq}\) and one non-abelian group.
    \end{enumerate}
\end{snippettheorem}

\begin{snippetproof}{p-q-groups-distinct-primes-product-order-theorem-proof}{p-q-groups-distinct-primes-product-order-theorem}{Classification of groups of order \(pq\)}
    Let \(P\) and \(Q\) be Sylow subgroups of orders \(p\) and \(q\).
    Sylow's theorem gives \(n_q\divides p\) and \(n_q\equiv1\pmod q\), so
    \(n_q=1\) and \(Q\normalsubgrp G\). Also \(n_p\) is \(1\) or \(q\),
    with the second possibility allowed only when \(q\equiv1\pmod p\).
    Since \(P\intersection Q=1\) and \(|PQ|=pq\), we have \(G=PQ\).

    If \(P\) is normal, \(G\) is the direct product
    \(C_p\cartesianprod C_q\groupisomorphic C_{pq}\). Thus this is the only
    possibility when \(q\not\equiv1\pmod p\).

    Suppose \(q\equiv1\pmod p\). The cyclic group
    \(\operatorname{Aut}(C_q)\) has order \(q-1\), so it contains an element
    of order \(p\). A non-trivial action
    \(C_p\fromto\operatorname{Aut}(C_q)\) therefore yields a non-abelian
    semidirect product \(C_p\ltimes C_q\). Conversely, every non-abelian
    group of order \(pq\) has this form. All non-trivial actions have the
    same image, the unique subgroup of order \(p\) in
    \(\operatorname{Aut}(C_q)\), and differ only by precomposition with an
    automorphism of \(C_p\). The semidirect-product isomorphism criterion
    therefore shows that they all give the same group up to isomorphism.
\end{snippetproof}

\section{Simple groups}

\begin{snippetdefinition}{simple-group-definition}{Simple group}
    A non-trivial \group \(G\) is \emph{simple} if its only
    \normalsubgrptext[normal subgroups] are \(1\) and \(G\).
\end{snippetdefinition}

\begin{snippetcorollary}{unique-proper-sylow-implies-not-simple}{A Sylow test for non-simplicity}
    Let \(G\) be a finite \group. If \(G\) has a unique non-trivial proper
    Sylow subgroup, then \(G\) is not simple.
\end{snippetcorollary}

\begin{snippetproof}{unique-proper-sylow-implies-not-simple-proof}{unique-proper-sylow-implies-not-simple}{A Sylow test for non-simplicity}
    A unique Sylow subgroup is normal. If it is non-trivial and proper, it is
    a forbidden normal subgroup in a simple group.
\end{snippetproof}

\begin{snippetexample}{group-order-280-not-simple}{Groups of order \(280\) are not simple}
    Let \(|G|=280=2^3\cdot5\cdot7\). Sylow's theorem gives
    \(n_7\in\{1,8\}\) and \(n_5\in\{1,56\}\). If either number is \(1\),
    then \(G\) is not simple. Otherwise, the Sylow subgroups of orders \(7\)
    and \(5\) contribute respectively \(8(7-1)=48\) and
    \(56(5-1)=224\) distinct non-identity elements. Together with the
    identity, these account for \(273\) elements, leaving exactly seven
    non-identity elements. A Sylow \(2\)-subgroup has order \(8\), so it must
    consist of the identity and those seven remaining elements. It is unique,
    hence normal, and \(G\) is not simple.
\end{snippetexample}

\end{document}
